% 318_Masseableitung_Geltungsbereich_En_ch.tex
% Scope of Mass Derivations in FFGFT:
% Rest Masses, Hadron Dynamics, and Open Bridges
% Doc. 318, August 2026


\section*{List of symbols}

\begin{center}
\begin{tabular}{lll}
\toprule
Symbol & Meaning & Value / unit \\
\midrule
$M_Z$ & Mass of the $Z$ boson; reference scale for $\alpha_s$ & $91.19$~GeV \\
$\Lambda_{\mathrm{QCD}}$ & QCD confinement scale & $\approx 200$--$217$~MeV \\
$\alpha_s(\mu)$ & Strong coupling at scale $\mu$ & $\alpha_s(M_Z) \approx 0.118$; $\alpha_s(m_\tau) \approx 0.33$ \\
$D = N_c$ & Number of colour degrees of freedom (Doc.~160) & $D = 3$ \\
$y_e$ & Yukawa coupling of the electron & $m_e/v \approx 2.08 \cdot 10^{-6}$ \\
$v$ & Higgs VEV & $246$~GeV \\
$m_p, m_n$ & Proton / neutron mass & $938.3$ / $939.6$~MeV \\
RGE & Renormalisation group equation & describes $\mu$-dependence of $\alpha_s$ \\
$K_{\text{frak}}$ & Fractal interface factor (rolled-up $\to$ SI) & $1-100\xi \approx 0.987$ \\
$m^{\text{bare}}$ / $m^{\text{SI}}$ & Geometric / measured mass & nat.~units / MeV \\
$\xi$ & FFGFT geometry parameter & $4/30000 \approx 1.333 \cdot 10^{-4}$ \\
\bottomrule
\end{tabular}
\end{center}

\section*{Summary}

The FFGFT derives particle masses from the geometry parameter
$\xi = 4/30000$ on the compact 4D torus $T^4/\mathbb{Z}_3$. This scope
is precise for rest masses from the Higgs-Yukawa sector (leptons, quark
rest masses). Hadron masses such as the proton mass contain a dominant
non-perturbative QCD contribution that requires its own framework. This
document draws the boundary explicitly and grounds it physically in the
relativistic energy relation $E^2 = (m_0 c^2)^2 + (pc)^2$ and the
picture of massless gluons inside the proton.

\begin{keyresult}[Status markers]
\textbf{[K]} FFGFT derives rest masses (leptons, quark Yukawa) from
$\xi$-geometry.\quad \textbf{Hadron masses} have additionally a
non-perturbative QCD contribution: \textbf{open bridge}.
\end{keyresult}

\section{Starting point: the complete energy relation}

The popular $E = m_0 c^2$ is an approximation. It follows from
linearising the Lorentz factor
\begin{equation}
  \gamma = \left(1 - \frac{v^2}{c^2}\right)^{-1/2}
  \approx 1 + \frac{1}{2}\frac{v^2}{c^2}
  \qquad (v \ll c)
\end{equation}
and holds only at small velocities. The complete, exact relation is:
\begin{equation}
  \boxed{E^2 = (m_0 c^2)^2 + (pc)^2}
  \label{eq:full_energy}
\end{equation}

This equation is the actual core -- not the simplified formula. The
second term $pc$ does not vanish when a particle is massless. For
massless particles ($m_0 = 0$):
\begin{equation}
  E = pc,
\end{equation}
so a massless particle moving at the speed of light still contributes
to the total energy -- and thus to the measurable mass of a system.

\section{The proton: mass as emergent property}

The proton mass $m_p \approx 938.3$~MeV is the clearest example of
this emergence:

\begin{table}[h]
\centering
\caption{Energy contributions to the proton mass}
\label{tab:proton_en}
\begin{tabular}{llll}
\toprule
Contribution & Value & Fraction & Mechanism \\
\midrule
$u$ and $d$ quark rest masses & $\approx 10$~MeV & $\approx 1\,\%$
  & Higgs-Yukawa coupling \\
Gluon kinetic energy & $\approx 928$~MeV & $\approx 99\,\%$
  & QCD confinement \\
\midrule
Proton mass & $938.3$~MeV & $100\,\%$ & \\
\bottomrule
\end{tabular}
\end{table}

The proton consists of three quarks ($uud$), whose rest masses together
amount to about $10$~MeV. The remainder -- $928$~MeV, around $99\,\%$
-- comes from gluons, the carriers of the strong force. Gluons are
massless ($m_0 = 0$). They move at the speed of light. By
Eq.~(\ref{eq:full_energy}) they nevertheless contribute to the total
energy of the system through their momentum.

\subsection{Timelessness inside the proton}

For massless particles moving at the speed of light, time dilation is
maximal: the proper time of these particles is zero. Gluons inside the
proton exist in a state where no time passes. For them, the spatial
extent of the world is zero -- length contraction is maximal.

And yet: these timeless, massless constituents account for the dominant
part of what is measured externally as the proton mass. Mass is a
property at the surface of a system, not an intrinsic property of its
components.

\textbf{[S]} This observation connects directly to the FFGFT fundamental
relation. $\tilde{T} \cdot m = 1$ states: intrinsic time and mass are
dual. The interior of the proton is precisely the timeless level -- and
this constitutes the majority of the externally apparent mass. Rest mass
arises through projection from the timeless to the time-bearing domain.
The proton illustrates $\tilde{T} \cdot m = 1$ inadvertently but
precisely. This connection is conceptually robust but not yet a formal
derivation; it is noted for a future document.

\section{Two sectors -- two mechanisms}

The ratio $m_p/m_e \approx 1836$ connects two physically distinct
sectors:

\begin{center}
\begin{tabular}{lll}
\toprule
Quantity & Sector & Scale \\
\midrule
$m_e, m_\mu, m_\tau$ & Higgs-Yukawa & $\xi^{p_i} \cdot r_i \cdot v$ \\
$m_{\nu_e}, m_{\nu_\mu}, m_{\nu_\tau}$ & Higgs-Yukawa (suppressed) & $\xi^{q_i} \cdot v$ \\
$m_u, m_d, m_s$ (rest) & Higgs-Yukawa & (open, Doc.~006) \\
$m_p, m_n$ (total) & QCD confinement + Higgs-Yukawa
  & $\Lambda_{\mathrm{QCD}} + m_{\mathrm{quarks}}$ \\
\bottomrule
\end{tabular}
\end{center}

\textbf{[B]} Proton mass and electron mass arise through entirely
different physics. $m_e$ is a rest mass from the Higgs-Yukawa sector;
$m_p$ consists to $99\,\%$ of QCD confinement energy. A ratio of these
two quantities connects two independent scales ($\Lambda_{\mathrm{QCD}}$
and $y_e \cdot v$). There is no structural reason why the geometry of a
single parameter $\xi$ should simultaneously fix both scales.

\textbf{Consequence:} $m_p/m_e$ is not a derivation target of the
current FFGFT corpus. An apparently exact five-digit agreement would be
a warning signal, not a success criterion -- because QCD contributions
carry percentage-level theory uncertainties.

\section{Neutrinos: the other limiting case}

In contrast to the proton, neutrinos are the opposite limit: extremely
small rest masses lying entirely within the Higgs-Yukawa sector.

\textbf{[S]} In FFGFT, neutrinos are treated as ``nearly massless
photons'' (Doc.~007): they carry the same $(n_\theta, n_\phi)$ winding
structure as the charged leptons, but with double $\xi$-suppression.
This yields rest masses in the sub-eV range, consistent with
cosmological bounds ($\sum m_\nu < 0.12$~eV, Planck 2018).

\begin{table}[h]
\centering
\caption{Neutrino quantum numbers and mass suppression (Doc.~007)}
\label{tab:neutrinos_en}
\begin{tabular}{lccccc}
\toprule
Neutrino & $n_\theta$ & $n_\phi$ & $p_i$ & Factor & Status \\
\midrule
$\nu_e$    & 3 & 2 & $5/2$ & $\xi^{5/2}$ & \textbf{[S]} \\
$\nu_\mu$  & 5 & 4 & $3$   & $\xi^{3}$   & \textbf{[S]} \\
$\nu_\tau$ & 9 & 5 & $25/9$& $\xi^{25/9}$& \textbf{[S]} \\
\bottomrule
\end{tabular}
\end{table}

Neutrinos demonstrate that the Higgs-Yukawa sector admits unlimited
suppression through powers of $\xi$. The proton lies outside this
sector.


\section{Rolled up and unrolled: the structural reason for the two-sector separation}

Doc.~311 distinguishes three ways in which the fourth spatial direction
of the 4D torus $T^4$ becomes the three-dimensional world:

\begin{center}
\begin{tabular}{lll}
\toprule
Way & Mechanism & Problem \\
\midrule
Compactification (Kaluza-Klein) & fourth direction becomes a circle, windings = modes
  & radius $R$ is a free modulus \\
Projection (cut-and-project) & 4D lattice $\to$ 3D cut
  & loses lattice invariants \\
Fourth is not space & carries mass, time, internal index
  & without compactness no mode tower \\
\bottomrule
\end{tabular}
\end{center}

\textbf{[K]} FFGFT combines ways 1 and 3, bound together by
$\tilde{T} \cdot m = 1$: the fourth direction is \emph{rolled up}
(compact) \emph{and} \emph{bound} (it carries mass). The radius $R$ is
therefore not free -- it is the mass. This is the only case in which a
scale arises without bringing in an additional parameter.

\subsection{Rolled up: rest masses from the Yukawa sector}

\textbf{[K]} Quark and lepton rest masses arise from \emph{rolled-up}
modes: winding numbers $(n_\theta, n_\phi)$ on the compact torus
generate quantised masses $m_i = r_i \cdot \xi^{p_i} \cdot v$
(Docs.~006, 317). Integer winding numbers are forced by the compactness
of the torus. This is the geometric origin of the Yukawa couplings: not
free parameters, but topological invariants.

\subsection{Unrolled: gluons as dynamic fields in 3D space}

\textbf{[S]} Gluons are massless vector bosons propagating in the
unrolled, non-compact 3D space. They are not winding modes on the torus
but dynamic degrees of freedom of the gauge field in the unfolded space.
Their energy does not come from $\xi$-windings but from motion and QCD
confinement.

\textbf{The proton in this language:}
\begin{center}
\begin{tabular}{lll}
\toprule
Contribution & Origin & Sector \\
\midrule
Quark rest masses ($\sim 1\,\%$) & rolled-up modes, Yukawa
  & FFGFT, \textbf{[K]} \\
Gluon kinetic energy ($\sim 99\,\%$) & unrolled modes, QCD dynamics
  & open bridge, \textbf{[S]} \\
\bottomrule
\end{tabular}
\end{center}

Hadron masses have contributions from \emph{both} regimes -- that is the
structural reason why a unified $\xi$-derivation does not directly apply.
The boundary between rolled-up and unrolled is the boundary of the scope.

\section{The fractal correction $K_{\text{frak}}$: interface factor in the transition rolled-up $\to$ unrolled}

\textbf{[K]} $K_{\text{frak}}$ is the factor that appears in the
transition from the geometrically rolled-up (bare) description to the
measured (unrolled, SI) description (Docs.~012, 133):
\begin{equation}
  m_i^{\text{SI}} = K_{\text{frak}} \cdot C_{\text{conv}} \cdot m_i^{\text{bare}},
  \qquad K_{\text{frak}} = 1 - 100\xi \approx 0.9867
\end{equation}
The physical justification is the fractal porosity of spacetime at
Planck scales ($D_{\text{frak}} = 2.94$, Doc.~003), which reduces the
effective interaction strength relative to a smooth space.

\subsection{The key point: $K_{\text{frak}}$ cancels in ratios}

\textbf{[K]} In mass ratios $K_{\text{frak}}$ drops out completely
(Doc.~012):
\begin{equation}
  \frac{m_\mu}{m_e} = \frac{K_{\text{frak}} \cdot m_\mu^{\text{bare}}}
  {K_{\text{frak}} \cdot m_e^{\text{bare}}} = \frac{m_\mu^{\text{bare}}}{m_e^{\text{bare}}}
\end{equation}

This is structurally significant: the mass ratios that FFGFT derives in
the lepton sector ($m_\mu/m_e$, $m_\tau/m_\mu$, Doc.~317) are purely
geometric and do \emph{not} require $K_{\text{frak}}$. $K_{\text{frak}}$
appears only when absolute SI values are needed -- for the gravitational
constant $G$, the fine-structure constant $\alpha$, and absolute mass
values in MeV or kg.

\subsection{Scope and limitation for hadrons}

$K_{\text{frak}}$ is the interface factor between two levels:
\begin{center}
\begin{tabular}{lll}
\toprule
Level & Description & $K_{\text{frak}}$ \\
\midrule
Rolled-up (bare) & $m^{\text{bare}}$, torus geometry, $\xi$-powers
  & not needed \\
Transition & $m^{\text{SI}} = K_{\text{frak}} \cdot C_{\text{conv}} \cdot m^{\text{bare}}$
  & \textbf{here} \\
Unrolled (SI) & measurable mass in MeV, kg & already included \\
\bottomrule
\end{tabular}
\end{center}

For the \textbf{hadron sector}, $K_{\text{frak}}$ would appear in the
transition $\Lambda_{\text{QCD}}^{\text{bare}} \to
\Lambda_{\text{QCD}}^{\text{SI}}$ -- but
$\Lambda_{\text{QCD}}^{\text{bare}}$ itself is not derived from $\xi$.
$K_{\text{frak}}$ as an interface factor does not resolve the underlying
problem: the $928$~MeV contribution is missing not because of a lacking
conversion, but because the underlying QCD mechanism is not yet treated
in the FFGFT corpus.

\subsection{A135: settability as a diagnostic}

\textbf{[K]} A135 states the settability criterion: what can be
norm-set to one was a property of the measurement system.
$K_{\text{frak}} = 1 - 100\xi$ cannot be set to~$1$ without changing
$\xi$ -- it is bound to something physical, not a pure measurement-system
artefact. For precision predictions Doc.~133 applies.

\section{Scope of the $\xi$ mass derivations}

\begin{table}[h]
\centering
\caption{Complete scope of mass derivations in FFGFT}
\label{tab:scope_en}
\begin{tabular}{llll}
\toprule
Particle & Dominant contribution & Status & Doc. \\
\midrule
$e, \mu, \tau$ & Higgs-Yukawa, rest mass
  & \textbf{[K]} & 006, 317 \\
$\nu_e, \nu_\mu, \nu_\tau$ & Higgs-Yukawa, $\xi^2$-suppressed
  & \textbf{[S]} & 007 \\
$u, d, s, c, b, t$ (rest) & Higgs-Yukawa
  & \textbf{[S]} open & 006 \\
$W, Z$ & Higgs mechanism
  & \textbf{[S]} scale correct & --- \\
Higgs $h$ & Higgs self-coupling
  & \textbf{[K]} $\xi_{\mathrm{Higgs}}$ & A-series \\
$p, n$ (total) & QCD + Higgs-Yukawa
  & \textbf{open bridge} & --- \\
$m_p/m_e$ & sector-mixed
  & \textbf{not claimed} & --- \\
\bottomrule
\end{tabular}
\end{table}

\section{Three scenarios for the QCD bridge}

Whether $\xi$ also captures the QCD sector is an open research question.
Three scenarios are conceivable:

\begin{enumerate}
  \item \textbf{$\xi$ fixes $\Lambda_{\mathrm{QCD}}$ via orbifold
    geometry.} The $T^4/\mathbb{Z}_3$ structure generates nine orbifold
    fixed points. Whether these force a natural QCD scale is open.
    A first step appears in Doc.~041:
    $\Lambda_{\mathrm{QCD}} = v \cdot \xi^{1/3} \approx 200$~MeV
    (experimental: $\sim 217$~MeV). This is a rough match but not a
    derivation -- the formula is listed there, not justified.
    \textbf{[S]}

  \item \textbf{Trace anomaly on compact spaces.} The renormalisation
    group for the strong coupling $\alpha_s(\mu)$ on compact tori
    generates a natural infrared scale. Khanna et al.\ (2014,
    arXiv:1409.1245) provide the mathematical framework; whether $\xi$
    enters there is unexplored. \textbf{[S]}

  \item \textbf{Hadron sector with SM inputs (Doc.~005) and
    geometric $\alpha_s$ derivation (Doc.~160).}
    Doc.~005 contains a complete proton calculation
    ($0.938100$~GeV, deviation $0.02\,\%$) using the extended fractal
    formula. It uses $\alpha_s = 0.118$ and
    $\Lambda_{\mathrm{QCD}} = 0.217$~GeV as Standard Model inputs.
    The claim ``without free parameters'' means: no additional
    parameters beyond the SM inputs.

    Essential precision:
    Doc.~160 (lepton lifetimes) derives $\alpha_s$ geometrically:
    \begin{equation}
      \alpha_s(m_\tau) = D \cdot \xi^{1/(D+1)} = 3 \cdot \xi^{1/4}
      \approx 0.322 \quad (\text{exp: }0.330,\; -2.3\,\%)
    \end{equation}
    $D = 3$ colour degrees of freedom, $D+1 = 4$ spacetime dimensions.
    This is $\alpha_s$ at scale $m_\tau$, not at $M_Z$.
    The strong coupling runs with energy: $\alpha_s(m_\tau)
    \approx 0.33$ and $\alpha_s(M_Z) \approx 0.118$ are
    different values of the \emph{same} running coupling.
    The FFGFT derivation targets the correct scale (Doc.~160)
    and uses $0.118$ for the proton calculation in Doc.~005
    because $M_Z$ is the relevant scale there. \textbf{[K]}

    Open point: the complete RGE bridge from
    $\alpha_s(m_\tau)$ to $\alpha_s(M_Z)$ within FFGFT
    is not yet carried out. \textbf{[S]}
\end{enumerate}

Scenario~3 is the current position. Scenarios~1 and~2 are marked
research questions. The approach in Doc.~041 ($\Lambda_{\mathrm{QCD}}
= v \cdot \xi^{1/3}$) is a first step towards Scenario~1 but not yet
a complete derivation.

\section{Relation to the corpus}

\subsection{Structural classification of the early documents}

\textbf{[K]} The early documents (Docs.~001--041) primarily worked out
the \emph{structural connections} between $\xi$ and the physical
constants -- not precision derivations. The formulas in these documents
are order-of-magnitude approaches: they show \emph{that} and \emph{how}
a coupling constant depends on $\xi$, not necessarily the exact
prefactor.

\begin{center}
\begin{tabular}{lll}
\toprule
Document & Content & Character \\
\midrule
Doc.~041 & $\alpha_s = \xi^{-1/3}$, $\alpha_W = \xi^{1/2}$ (nat.~units)
  & structural sketch \textbf{[S]} \\
Doc.~041 & $\Lambda_{\mathrm{QCD}} = v \cdot \xi^{1/3} \approx 200$~MeV
  & order of magnitude \textbf{[S]} \\
Doc.~160 & $\alpha_s(m_\tau) = 3\xi^{1/4} \approx 0.322$
  & precision derivation \textbf{[K]} \\
Doc.~011 & $\alpha$ from torus geometry
  & precision derivation \textbf{[K]} \\
Doc.~006 & lepton masses $m_i = r_i\xi^{p_i}v$
  & precision derivation \textbf{[K]} \\
\bottomrule
\end{tabular}
\end{center}

In particular: Doc.~041 lists $\alpha_s \approx 0.118$ (experiment)
and $\xi^{-1/3} = 9.65$ (natural units) side by side without unit
conversion. This is not an error in the physics but a missing bridge --
the conversion from natural to SI units is absent there. Doc.~160
provides the cleaner derivation.

\subsection{Corpus references}

\begin{itemize}
  \item \textbf{Doc.~005} (tm-extension): proton calculation with
    extended fractal formula; $\alpha_s = 0.118$ and
    $\Lambda_{\mathrm{QCD}} = 0.217$~GeV as SM inputs (not $\xi$-derivations).
  \item \textbf{Doc.~006} (particle masses): lepton rest masses [K];
    hadron masses not treated.
  \item \textbf{Doc.~007} (neutrinos): double $\xi$-suppression,
    quantum number assignment, cosmological bounds.
  \item \textbf{Doc.~041} (parameter derivation): structural connections
    of all couplings; order-of-magnitude estimates, not SI precision.
  \item \textbf{Doc.~160} (lepton lifetimes): geometric derivation
    $\alpha_s(m_\tau) = 3\xi^{1/4}$, precision $-2.3\,\%$ [K].
  \item \textbf{Doc.~317} (KSAU-FFGFT): complete quantum number table.
  \item \textbf{A-series}: $\xi_{\mathrm{Higgs}}$, Koide relation.
\end{itemize}

\section{Entry in Doc.~190 (R76)}

This clarification is declared as register entry R76 in Doc.~190:
scope of mass derivations explicitly delimited; hadron sector documented
as open bridge; $m_p/m_e$ not claimed as a derivation target of the
current corpus; three scenarios for the QCD bridge marked as research
questions.
